% !TeX root = ../papers/nc-p3.tex
\section{相補逆行列の半単純スペクトルと四元群印を持たない孤立型}
\label{sec:split-bell-locus}
前節の $P$ の第一切断を四次元空間の間の可逆行列 $M$ と見る。
各側は二つの二次元空間の積なので行列式二次形式を持つ。その形式による随伴を $M^\dagger$ とする。
ここでいう半単純とは $M^\dagger M$ が対角化可能であることを指す。
\begin{theorem}\label{thm:split-bell-locus}
半単純な $P$ は以下の射影四数 $x=(a,b,c,d)$ で、座標置換と偶数個の符号反転を除いて分類される。
\[
  f_0=abcd,\qquad
  -16 f_1=(a-b-c-d)(a-b+c+d)(a+b-c+d)(a+b+c-d).
\]
$f_0f_1\ne0$ を課す。前節の移送集合が非空となる $P$ は、以下とその付け替えで尽くされる。
\begin{enumerate}
\item 曲面 $ab+cd=0$。
\item 直線 $a=b,c=d$。
\item $a=1$, $(d^2+1)^3=4d^3$,
$2c=d^3+d^2+3d-1$, $2b=d^5-d^4+3d^3-7d^2+6d-2$。
この六点は付け替えの下でちょうど二つの軌道を与える。
\end{enumerate}
曲面の一般点の上には分解付きで四つの添字同型類がある。第三型では二軌道の各々に
ちょうど一つあり、分解を忘れると両者は同型である。この一つの標準AS窓は最小周期1で、
四位置の射影安定群は $\mathbb G_m$ 自身である。従って四元群印を持つ窓のどれとも同型でない。
孤立は本定理の $P$ の分類についての主張である。
\end{theorem}
\begin{proof}[標準形と九十六通りの比較]
$\operatorname{SL}_2\times\operatorname{SL}_2$ の四次元像は行列式形式の特殊直交群である
（核は同時の $\pm I$、像は同次元の連結閉部分群）。
$M^\dagger M$ の異なる固有空間は互いに直交し、各空間は非退化である。
各々の正規直交基底を右側に取り、その $M$ 像を固有値の平方根で割ると左側の
正規直交基底を得る。向きを揃えると、対角の四係数の自由度は順序と偶符号変更および共通倍率だけである。
これは固有値が重なっても同じで、二乗の多重集合と四係数の積が残る向きを決める。
通常の二進基底に戻せば非零になり得る四対は
$0000,0011,0101,0110$ とその全ビット反転で、各対の値は
\[
             (a+b)/2,\ (a-b)/2,\ (c+d)/2,\ (c-d)/2
\]
である。この規約は二側のBell基底 $00\pm11,01\pm10$ を使うもので、
二次形式の符号に合わせた単位因子を両側で逆に取れる。
四成分の一斉符号は射影的に消えるので、有限群 $W$ の位数は96である。

両切断の行列式は $f_0,f_1$ であり、逆転置の射影Bell四数はそれぞれ
$\nabla f_0,\nabla f_1$ である。これは二つの二行列の逆転置を直接取った式で、
二次元双対を同一視する交代形式四つもこの表示を保つ。
二重互換 $\tau$ は四対の値を変えない。従って比較は $g=(\pi,e)\in W$ に対する三条件
\[
  x_i e_i\,\partial_{\pi(i)}f_1=x_0 e_0\,\partial_{\pi(0)}f_1
                    \qquad(i=1,2,3)
\]
に等しい。各射影勾配は $W$ 同変である。実際、四文字の局所基底変換に対し
両切断の逆転置はともに双対に移り、交代形式による射影同定は同変である。
従って共役な $g$ の解は $W$ で移る。

$a=1$ として十共役類の各代表で三式の極小素成分を求め、$f_0f_1=0$ に入る成分を除く。
\path{computations/verify_split_bell_locus.py} は四つの非零座標と四一次因子を個別に調べる。
代表の順に共役類の大きさは $3,1,12,12,32,6,3,3,12,12$、残る有理素成分数は
$0,6,0,2,2,1,0,3,0,3$ である。各素成分で三曲面の積の式または六直線の一つの二式が
消えることを確認し、残る一成分では消去イデアルが第三型の三式と一致する。
逆向きには恒等変換で第二型、$(01)(23)$ で第一型の三式が恒等的に消える。
第五代表の六次式は既約で、三曲面式の積の剰余は $-4$ である。
六点では $\sum x_i^2\ne0$ かつ平方は全て相異なるので、$[x]$ を保つ $W$ は自明である。
第一逆写像は対合だからその像の $W$ 安定群も自明で、六点それぞれの比較 $g$ は一意である。
第五代表の32共役を動かして得る $32\cdot6$ 点は従って相異なり、二つの96点軌道となる。
\end{proof}

\begin{proof}[曲面の一般ファイバーと孤立窓]
平方が相異なり、$\sum x_i^2\ne0$ なら $D$ の安定群 $H$ はちょうど四元群である。
左右の直交行列は共通の四固有線上の偶符号変更となり、全符号を除くと四つ残る。
元の各二次元因子の射影線では、全位置同時のビット反転と全位置同時の
$\operatorname{diag}(1,-1)$ が生成する。$g$ のベル四線の置換が $H$ の三つの非自明元に
誘導する作用は、四線の三つの二対分割への置換である。各偶符号はその正負で四線を二分するためである。
$\tau$ は上の二生成元をそれぞれ保つ。従って曲面の一般点で $g=(01)(23)$ は
$H$ に自明に作用し、前節の $Hv$ の共役類は四個ある。
六次式では第五代表の置換が三サイクルで $H$ の非自明三元を巡回する。
このとき $1-\operatorname{Ad}(v)$ は $H$ 上全単射なので $Hv$ は単一共役類である。

分解を交換すると第一交代成分は旧 $Q$ になり、最初の四対応で座標を揃えると
そのBell四数は $[1/a:1/b:1/c:1/d]$ である。上記検算は六次式を法として
この四数を元に戻す $W$ 元が一つもないことを確認する。交換は対合であり、
曲面・直線の付け替えの和集合は逆数で保たれるので二つの孤立軌道を交換する。
一位置移動の新しい交代成分は $Q$ の一巡回である。同じ六次剰余体で、
そのBell四数を旧 $P$ に戻す $W$ 元の存在も検算する。移送の共役類は一つなので
分解付きでも最小周期1になる。

最後に孤立窓の多安定性を確かめる。複素数上で、各BellベクトルをEuclidノルム1にし、
行列式形式に対する符号が負のものに虚数単位を掛けると、対称形式とHermite形式の
双方について正規直交となる。行列式1の実符号置換はこの基底の実特殊直交群に属し、
二つの二次元因子でユニタリ行列に持ち上がる。実際この実群は
$\operatorname{SU}(2)\times\operatorname{SU}(2)$ の像である
（コンパクトな連結六次元群で、二つの形式を保つ）。右左で必要な置換・偶符号も
この持上げを用いる。第一の四対応と第二の四対応を射影ユニタリに選べるので、
四つの境界成分は最大階数の二文字ユニタリ形式二つの積である。
Bell型の $P,Q$ でも、各一文字Gram行列は偶奇の台と全反転によりスカラーである。
六つの二重次数成分は各一文字の同じ半分で異なる相手台を持つ。
各高位半分の実倍率を動かすと、六つの二要素集合の重みの凸包の相対内部に
均一周辺和の点があるので、全位置で両半分のGram値を一致させられる
（節\ref{sec:klein-polystable-limits}のノルム最小化）。
よって全基底群で多安定で、射影安定群は簡約群となる。

指定トーラス $T$ の中心化群を $T$ で割ったものは前節の構成で $H$ のうち
$v$ と可換な群に入り、ここでは自明である。$T$ の反転は分解の交換に当たり、
二軌道が印付きで相異なるので安定群内で起こらない。従って $N(T)=T$。
$T$ は極大トーラスで、連結簡約群の非自明な根は非自明なWeyl反射を与えるから
安定群の恒等成分は $T$ である。どの連結成分にも $T$ の正規化元があるため全群も $T$ となる。
\end{proof}
